% ============================================================================
%  TCN Configuration Comparison Report
%  Ibrahim Hanafy — NU Research Internship 2026
% ============================================================================
\documentclass[11pt, a4paper]{article}

% ── Packages ────────────────────────────────────────────────────────────────
\usepackage[margin=2.2cm]{geometry}
\usepackage{graphicx}
\usepackage{booktabs}
\usepackage{amsmath}
\usepackage{hyperref}
\usepackage{xcolor}
\usepackage{caption}
\usepackage{subcaption}
\usepackage{float}
\usepackage[sorting=none]{biblatex}
\usepackage{enumitem}

% ── Colors ──────────────────────────────────────────────────────────────────
\definecolor{nublue}{HTML}{1565C0}
\definecolor{accent}{HTML}{FF7043}

\hypersetup{
    colorlinks=true,
    linkcolor=nublue,
    citecolor=nublue,
    urlcolor=nublue,
}

% ── Title ───────────────────────────────────────────────────────────────────
\title{%
    \textcolor{nublue}{\rule{\linewidth}{1.2pt}}\\[6pt]
    \textbf{TCN Multi-Output ECG Forecasting}\\[4pt]
    Configuration Comparison Report\\[6pt]
    \textcolor{nublue}{\rule{\linewidth}{1.2pt}}
}
\author{Ibrahim Hanafy\\[2pt]
    \small Nile University --- Research Internship 2026\\
    \small \texttt{s-ibrahim.hanafy@zewailcity.edu.eg}}
\date{August 12, 2026}

\begin{document}
\maketitle

% ════════════════════════════════════════════════════════════════════════════
\section{Objective}
% ════════════════════════════════════════════════════════════════════════════

This report compares two Temporal Convolutional Network (TCN) configurations for multi-step ECG forecasting on the MIT-BIH Arrhythmia Database. The goal is to demonstrate that our modified configuration outperforms the original paper architecture, thereby justifying systematic hyperparameter optimization as the next research phase.

% ════════════════════════════════════════════════════════════════════════════
\section{Experimental Setup}
% ════════════════════════════════════════════════════════════════════════════

\subsection{Dataset}
\begin{itemize}[nosep]
    \item \textbf{Source:} MIT-BIH Arrhythmia Database (21 patients, excluding records 111 \& 118)
    \item \textbf{Strategy:} Multi-Output (MIMO) --- single model predicts all future steps simultaneously
    \item \textbf{Horizons:} $H \in \{1, 10, 15, 20\}$ steps ahead
\end{itemize}

\subsection{Configuration Comparison}

\begin{table}[H]
\centering
\caption{Hyperparameter configurations compared.}
\label{tab:configs}
\begin{tabular}{lcc}
\toprule
\textbf{Hyperparameter} & \textbf{Paper Config} & \textbf{Our Config} \\
\midrule
Residual Blocks     & 1            & \textbf{3} (dilations: 1, 2, 4) \\
Filters             & 256          & \textbf{64} \\
Batch Size          & 128          & \textbf{256} \\
Kernel Size         & 3            & 3 \\
Dropout Rate        & 0.1          & 0.1 \\
Learning Rate       & $10^{-3}$    & $10^{-3}$ \\
Epochs              & 200          & 200 \\
Early Stopping      & patience 50  & patience 50 \\
Optimizer           & Adam         & Adam \\
\bottomrule
\end{tabular}
\end{table}

\noindent\textbf{Key difference:} Our configuration trades wide filters (256) for a deeper architecture (3 residual blocks with 64 filters each), increasing the effective receptive field through stacked dilations while using $4\times$ fewer parameters per layer.

% ════════════════════════════════════════════════════════════════════════════
\section{Results}
% ════════════════════════════════════════════════════════════════════════════

\subsection{Aggregate Performance}

\begin{table}[H]
\centering
\caption{Mean metrics across all 21 patients. \textbf{Bold} = better.}
\label{tab:results}
\begin{tabular}{c cc cc cc}
\toprule
& \multicolumn{2}{c}{\textbf{RMSE} $\downarrow$}
& \multicolumn{2}{c}{\textbf{MAE} $\downarrow$}
& \multicolumn{2}{c}{\textbf{R\textsuperscript{2}} $\uparrow$} \\
\cmidrule(lr){2-3} \cmidrule(lr){4-5} \cmidrule(lr){6-7}
\textbf{H} & Ours & Paper & Ours & Paper & Ours & Paper \\
\midrule
1  & \textbf{0.0052} & 0.0061 & \textbf{0.0036} & 0.0037 & \textbf{0.9970} & 0.9952 \\
10 & \textbf{0.0304} & 0.0344 & \textbf{0.0134} & 0.0148 & \textbf{0.8991} & 0.8776 \\
15 & \textbf{0.0433} & 0.0471 & \textbf{0.0185} & 0.0200 & \textbf{0.8108} & 0.7847 \\
20 & \textbf{0.0539} & 0.0571 & \textbf{0.0236} & 0.0253 & \textbf{0.7157} & 0.6894 \\
\bottomrule
\end{tabular}
\end{table}

\noindent Our configuration achieves lower RMSE, lower MAE, and higher R\textsuperscript{2} at \emph{every} forecast horizon.

\subsection{Win Rate Analysis}

\begin{table}[H]
\centering
\caption{Number of patients (out of 21) where our config outperforms.}
\label{tab:wins}
\begin{tabular}{ccc}
\toprule
\textbf{Horizon} & \textbf{RMSE Wins} & \textbf{R\textsuperscript{2} Wins} \\
\midrule
1  & 12/21 (57\%) & 12/21 (57\%) \\
10 & 19/21 (90\%) & 19/21 (90\%) \\
15 & 20/21 (95\%) & 20/21 (95\%) \\
20 & 19/21 (90\%) & 19/21 (90\%) \\
\bottomrule
\end{tabular}
\end{table}

\noindent At longer horizons ($H \geq 10$), our config wins on 90--95\% of patients, confirming that deeper dilated architectures generalize better for multi-step forecasting.

\subsection{Training Efficiency}

\begin{table}[H]
\centering
\caption{Mean training time per patient (seconds).}
\label{tab:time}
\begin{tabular}{ccc}
\toprule
\textbf{Horizon} & \textbf{Ours (s)} & \textbf{Paper (s)} \\
\midrule
1  & \textbf{169.8} & 226.6 \\
10 & \textbf{159.5} & 207.7 \\
15 & \textbf{160.4} & 202.0 \\
20 & \textbf{156.4} & 188.6 \\
\bottomrule
\end{tabular}
\end{table}

\noindent Our config trains 17--25\% faster due to $4\times$ fewer filters (64 vs.\ 256), despite having 3 residual blocks instead of 1.

% ════════════════════════════════════════════════════════════════════════════
\section{Visual Analysis}
% ════════════════════════════════════════════════════════════════════════════

\begin{figure}[H]
\centering
\begin{subfigure}[b]{0.48\textwidth}
    \includegraphics[width=\textwidth]{Figures/report_plots/bar_rmse.png}
    \caption{Mean RMSE by horizon}
\end{subfigure}
\hfill
\begin{subfigure}[b]{0.48\textwidth}
    \includegraphics[width=\textwidth]{Figures/report_plots/bar_r2.png}
    \caption{Mean R\textsuperscript{2} by horizon}
\end{subfigure}
\caption{Our configuration (blue) consistently outperforms the paper configuration (orange).}
\label{fig:bars}
\end{figure}

\begin{figure}[H]
\centering
\begin{subfigure}[b]{0.48\textwidth}
    \includegraphics[width=\textwidth]{Figures/report_plots/bar_mae.png}
    \caption{Mean MAE by horizon}
\end{subfigure}
\hfill
\begin{subfigure}[b]{0.48\textwidth}
    \includegraphics[width=\textwidth]{Figures/report_plots/bar_time.png}
    \caption{Mean training time by horizon}
\end{subfigure}
\caption{Our config achieves lower MAE while training faster.}
\label{fig:mae_time}
\end{figure}

\begin{figure}[H]
\centering
\includegraphics[width=0.95\textwidth]{Figures/report_plots/box_rmse.png}
\caption{RMSE distribution across patients for each horizon. Our config shows tighter distributions (lower variance) and lower medians.}
\label{fig:box_rmse}
\end{figure}

\begin{figure}[H]
\centering
\includegraphics[width=0.95\textwidth]{Figures/report_plots/heatmap_r2_diff.png}
\caption{Per-patient $\Delta$R\textsuperscript{2} = Ours $-$ Paper. Blue bars indicate patients where our config achieves higher R\textsuperscript{2}. At $H \geq 10$, nearly all patients benefit.}
\label{fig:heatmap}
\end{figure}

% ════════════════════════════════════════════════════════════════════════════
\section{Key Findings}
% ════════════════════════════════════════════════════════════════════════════

\begin{enumerate}[nosep]
    \item \textbf{Depth over width:} 3 residual blocks with 64 filters outperforms 1 block with 256 filters across all metrics and horizons.
    \item \textbf{Stronger at longer horizons:} The advantage grows from 57\% win rate at $H=1$ to 90--95\% at $H \geq 10$, indicating deeper dilated convolutions better capture long-range temporal dependencies.
    \item \textbf{Computational efficiency:} Our config trains 17--25\% faster due to fewer parameters, making it more practical for large-scale patient-wise experiments.
    \item \textbf{Consistent improvement:} Lower variance in RMSE distributions suggests our architecture generalizes more reliably across different patient morphologies.
\end{enumerate}

% ════════════════════════════════════════════════════════════════════════════
\section{Recommendation: Hyperparameter Optimization}
% ════════════════════════════════════════════════════════════════════════════

These results confirm that even a manual adjustment (depth vs.\ width) yields meaningful gains. This motivates a \textbf{systematic hyperparameter optimization} phase to explore the full configuration space:

\begin{itemize}[nosep]
    \item \textbf{Search space:} Number of blocks (2--6), filters (32--256), kernel size (2--7), dropout (0.05--0.3), learning rate ($10^{-4}$ to $10^{-2}$), batch size (64--512)
    \item \textbf{Methods:} Bayesian optimization (Optuna), random search, or Hyperband
    \item \textbf{Objective:} Minimize mean RMSE across horizons, with R\textsuperscript{2} as secondary metric
    \item \textbf{Validation:} Patient-wise cross-validation to ensure generalizability
\end{itemize}

\vspace{1em}
\noindent\textcolor{nublue}{\rule{\linewidth}{0.8pt}}\\[4pt]
\noindent\textit{Report generated August 12, 2026. Data from 21 MIT-BIH patients, 4 horizons, 2 TCN configurations.}

\end{document}
